\subsection{Measurement-Driven Motivation}
\label{sec:motivation}

We conduct a controlled ablation study to identify where the performance bottleneck lies in real-time 3DGS streaming.
Using Spark~2.0's streaming-LoD pipeline on a 40\,M-splat scene, we measure three workload scenarios under three delivery tiers (3 trials each), totaling 308+ trials.

\paragraph{Workload scenarios.}
We decompose the streaming pipeline into three isolation levels:
\begin{itemize}[leftmargin=*,nosep]
\item \textbf{Normal streaming}: the full pipeline---LoD chunks fetched via HTTP Range, decoded, uploaded to GPU, and rendered.
\item \textbf{Fetch-only}: network I/O and decode proceed but geometry is discarded before GPU upload, isolating network and parse cost.
\item \textbf{Pure render}: chunk fetching blocked during measurement; the renderer operates on pre-loaded geometry only.
\end{itemize}

\paragraph{Delivery tiers.}
(1)~A \emph{local baseline} where the asset server runs on the same machine (localhost/loopback reference path);
(2)~an \emph{origin server} hosting the .rad file over HTTP with Range support (measured chunk p95 completion = 499\,ms per 1.7\,MB chunk, $n{=}3$);
(3)~a \emph{remote server} serving the same file from a geographically distributed edge (measured chunk p95 = 15{,}872\,ms, $n{=}3$).
All trials use Chrome~147 on an NVIDIA RTX~4070 Laptop GPU at 1280$\times$720.

\begin{table}[t]
\centering
\caption{Workload isolation across three delivery tiers (single-user, 3 trials each). \emph{Chunks} = Range requests completed in the 40\,s measurement window. \emph{Chunk p95} = p95 completion time of a single ${\sim}$1.7\,MB Range request.}
\label{tab:motivation}
\small
\begin{tabular}{@{}llrrrr@{}}
\toprule
Delivery & Scenario & FPS & Chunks & Chunk p95 (ms) \\
\midrule
Local baseline & Normal streaming & 86.9 & 230 & 30 \\
Local baseline & Fetch-only & 129.2 & 140 & 161 \\
Local baseline & Pure render & 115.5 & 0 & --- \\
\midrule
Origin server & Normal streaming & 88.3 & 277 & 499 \\
Origin server & Fetch-only & 165.0 & 140 & 991 \\
Origin server & Pure render & 92.3 & 0 & --- \\
\midrule
Remote server & Normal streaming & 156.3 & 17 & 15{,}872 \\
Remote server & Fetch-only & 165.0 & 5 & 140{,}626 \\
Remote server & Pure render & 164.9 & 0 & --- \\
\bottomrule
\end{tabular}
\end{table}

\paragraph{Layer 1: Local baseline confirms the client pipeline is functional.}
Under local delivery (Table~\ref{tab:motivation}, top), normal streaming sustains 86.9\,fps while completing 230 chunk requests in 40\,s.
Pure render achieves 115.5\,fps on pre-loaded geometry.
The gap reflects decode and GPU upload cost---measurable but not the dominant factor.
Pure rendering is not the primary bottleneck under this setup.

\paragraph{Layer 2: Origin server exposes per-chunk latency; p95 increases across tiers.}
When chunks are served from an origin server ($n{=}3$ trials), normal streaming achieves 88.3\,fps with 277 chunks---comparable to local in throughput.
However, the p95 chunk completion time rises from 30\,ms to 499\,ms (17$\times$).
Each 1.7\,MB Range request now takes hundreds of milliseconds; the pipeline sustains aggregate throughput only because multiple requests overlap in flight.
Fetch-only reveals the network cost more directly: p95 rises to 991\,ms, and the same 140 chunks complete but with substantially higher per-request latency.
Cold-start time to first visible splat is 1{,}178\,ms ($n{=}3$; range 948--1{,}657\,ms).

\paragraph{Layer 3: Remote server amplifies the mismatch; ``fast but empty'' rendering emerges.}
Under remote delivery, normal streaming reports 156.3\,fps---\emph{paradoxically higher} than local or origin.
But only 17 chunks arrive in 40\,s (93\% fewer than local), with p95 chunk completion of 15{,}872\,ms.
The elevated frame rate is a ``fast but empty'' artifact: so few chunks arrive that decode and upload rarely execute, leaving the render loop uncontested---but the displayed scene has drastically lower geometric detail.
Remote chunk arrival drops sharply while p95 increases by 500$\times$ relative to local.
Cold-start first visible splat rises to 6{,}703\,ms (5.7$\times$ slower than origin).

Under trace-based sequential 4-user replay, remote delivery completes only 12 chunks per client (vs.\ 288 locally) with p95 of 21{,}141\,ms---because the server treats each Range request independently regardless of viewport overlap across users.

\paragraph{The fundamental mismatch.}
These results show that on-demand large Range fetching---the standard HTTP mechanism for streaming LoD data---is incompatible with interactive 3DGS quality requirements once delivery leaves the local machine.
At the origin server, per-chunk p95 is already 17$\times$ above local; at a remote server it is 500$\times$ above.
The server has no knowledge of which chunks the viewport will need next, which chunks are shared across users, or which source would deliver fastest.

This motivates \emph{Viewport-aware Range Coordination} (VRC): a system that performs viewport prediction to prefetch chunks before they are needed, assigns chunk priority based on visibility contribution, deduplicates overlapping requests across users, and schedules fetches across sources based on measured latency---without modifying the underlying HTTP Range servers.
